\documentclass[tikz,border=10pt]{standalone}
\usepackage{tikz}
\usetikzlibrary{arrows.meta, positioning, calc}

\begin{document}
% Custom colours sampled from the image
\definecolor{catborder}{HTML}{000000}
\definecolor{algborder}{HTML}{7A9DBC}
\definecolor{hlfill}{HTML}{D9E6F5}

\begin{tikzpicture}[
    % y=0.95cm, % <--- You can now adjust this freely!
    font=\sffamily,
    >=Latex,
    % Styles for the category boxes
    category/.style={
      rectangle, rounded corners=8pt, draw=catborder, thick,
      fill=white, minimum width=4.0cm, minimum height=1.2cm,
      align=center, font=\normalsize
    },
    % Styles for the algorithm boxes
    algo/.style={
      rectangle, rounded corners=4pt, draw=algborder, thick,
      fill=white, minimum width=3.4cm, minimum height=0.9cm,
      align=center, font=\normalsize
    },
    algohl/.style={
      algo, fill=hlfill
    },
    trunk/.style={thick},
    branch/.style={->, thick}
  ]

  % ---------------------------------------------------------
  % 1. CATEGORY NODES
  % ---------------------------------------------------------
  \node[category] (rl)  at (0.75, 0)     {RL Algorithms};

  \node[category] (mf)  at (-4.5, -2.2)  {Model-Free RL};
  \node[category] (mb)  at (6.0, -2.2)   {Model-Based RL};

  \node[category] (po)  at (-7.0, -4.4)  {Policy Optimisation};
  \node[category] (ql)  at (-2.0, -4.4)  {Q-Learning};

  \node[category] (ltm) at (3.5, -4.4)   {Learn the Model};
  \node[category] (gtm) at (8.5, -4.4)   {Given the Model};

  % ---------------------------------------------------------
  % 2. ALGORITHM NODES
  % ---------------------------------------------------------
  % Col 1: Under Policy Optimisation (Left)
  \node[algo]   (pg)   at (-10.0, -6.2)  {Policy Gradient};
  \node[algo]   (a2c)  at (-10.0, -7.8)  {A2C / A3C};
  \node[algohl] (ppo)  at (-10.0, -9.4)  {PPO};
  \node[algo]   (trpo) at (-10.0, -11.0) {TRPO};

  % Col 3: Shared between Policy Optimisation and Q-Learning
  \node[algo]   (ddpg) at (-4.5, -7.0)   {DDPG};
  \node[algohl] (td3)  at (-4.5, -8.6)   {TD3};
  \node[algohl] (sac)  at (-4.5, -10.2)  {SAC};

  % Col 5: Under Q-Learning (Right)
  \node[algo]   (dqn)  at (1.0, -6.2)    {DQN};
  \node[algo]   (c51)  at (1.0, -7.8)    {C51};
  \node[algo]   (qr)   at (1.0, -9.4)    {QR-DQN};
  \node[algo]   (her)  at (1.0, -11.0)   {HER};

  % Col 7: Under Learn the Model (Right)
  \node[algo]   (wm)   at (6.5, -6.2)    {World Models};
  \node[algo]   (i2a)  at (6.5, -7.8)    {I2A};
  \node[algo]   (mbmf) at (6.5, -9.4)    {MBMF};
  \node[algo]   (mbve) at (6.5, -11.0)   {MBVE};

  % Col 9: Under Given the Model (Right)
  \node[algo]   (az)   at (11.5, -7.0)   {AlphaZero};

  % ---------------------------------------------------------
  % 3. CONNECTIONS (Now 100% Dynamic)
  % ---------------------------------------------------------

  \tikzset{splitline/.style={->, thick, rounded corners=6pt}}

  % Top Level Splits
  % Calculates exactly the halfway point between row 1 and row 2
  \path (rl.south) -- (rl.south |- mf.north) coordinate[midway] (fork1);
  \draw[thick] (rl.south) -- (fork1);
  \draw[splitline] (fork1) -| (mf.north);
  \draw[splitline] (fork1) -| (mb.north);

  % Mid Level Splits
  \path (mf.south) -- (mf.south |- po.north) coordinate[midway] (fork2);
  \draw[thick] (mf.south) -- (fork2);
  \draw[splitline] (fork2) -| (po.north);
  \draw[splitline] (fork2) -| (ql.north);

  \path (mb.south) -- (mb.south |- ltm.north) coordinate[midway] (fork3);
  \draw[thick] (mb.south) -- (fork3);
  \draw[splitline] (fork3) -| (ltm.north);
  \draw[splitline] (fork3) -| (gtm.north);

  % --- Spines (Vertical Trunks) ---
  % Spines now draw automatically to the exact height of the lowest node they serve
  \draw[trunk] (po.south)  -- (po.south |- trpo.east);
  \draw[trunk] (ql.south)  -- (ql.south |- her.west);
  \draw[trunk] (ltm.south) -- (ltm.south |- mbve.west);
  \draw[trunk] (gtm.south) -- (gtm.south |- az.west);

  % --- Branches (Horizontal Connections) ---
  % They automatically trace from the spine to the side of the node

  % Policy Optimisation
  \draw[branch] (po.south |- pg.east)   -- (pg.east);
  \draw[branch] (po.south |- ddpg.west) -- (ddpg.west);
  \draw[branch] (po.south |- a2c.east)  -- (a2c.east);
  \draw[branch] (po.south |- td3.west)  -- (td3.west);
  \draw[branch] (po.south |- ppo.east)  -- (ppo.east);
  \draw[branch] (po.south |- sac.west)  -- (sac.west);
  \draw[branch] (po.south |- trpo.east) -- (trpo.east);

  % Q-Learning
  \draw[branch] (ql.south |- dqn.west)  -- (dqn.west);
  \draw[branch] (ql.south |- ddpg.east) -- (ddpg.east);
  \draw[branch] (ql.south |- c51.west)  -- (c51.west);
  \draw[branch] (ql.south |- td3.east)  -- (td3.east);
  \draw[branch] (ql.south |- qr.west)   -- (qr.west);
  \draw[branch] (ql.south |- sac.east)  -- (sac.east);
  \draw[branch] (ql.south |- her.west)  -- (her.west);

  % Learn the Model
  \draw[branch] (ltm.south |- wm.west)   -- (wm.west);
  \draw[branch] (ltm.south |- i2a.west)  -- (i2a.west);
  \draw[branch] (ltm.south |- mbmf.west) -- (mbmf.west);
  \draw[branch] (ltm.south |- mbve.west) -- (mbve.west);

  % Given the Model
  \draw[branch] (gtm.south |- az.west)   -- (az.west);

\end{tikzpicture}
\end{document}
